\documentclass[../../../include/open-logic-section]{subfiles}

\begin{document}

\olfileid{sth}{choice}{vitali}
\olsection{પરિશિષ્ટ: વિટાલીનો વિરોધાભાસ}

બાનાખ--ટાર્સ્કીની રચના સ્વીકાર્ય છે કે નહીં
તેનો સાચો ખ્યાલ મેળવવો હોય, તો તેનો
\emph{પુરાવો} જોવો જોઈએ. દુર્ભાગ્યે તેના
માટે અહીં રજૂ કરી શકીએ તે કરતાં ઘણું
વધુ બીજગણિત જોઈએ. જોકે નીચેના
પરિણામ પર ધ્યાન આપીને બાનાખ--ટાર્સ્કીના
પુરાવા પર પ્રકાશ પાડે એવી કેટલીક ટૂંકી
ટિપ્પણીઓ કરી શકીએ.\footnote{વધુ વિસ્તૃત
ચર્ચા માટે \cite{Weston2003} અથવા
\cite{Wagon2016} જુઓ.}

\begin{thm}[વિટાલીનો વિરોધાભાસ ($\ZFC$માં)]\ollabel{vitaliparadox}
કોઈ પણ વર્તુળને ગણનીય સંખ્યામાં ટુકડાઓમાં
વિભાજિત કરી શકાય છે; એ ટુકડાઓને
(પરિભ્રમણ અને સ્થાનાંતરણથી) ફરી ગોઠવી
તે વર્તુળની બે નકલો બનાવી શકાય છે.
\end{thm}

વિટાલીનો વિરોધાભાસ બાનાખ--ટાર્સ્કીના
વિરોધાભાસ કરતાં સાબિત કરવામાં ઘણો
સરળ છે. તેને ``વિટાલીનો વિરોધાભાસ''
કહ્યો છે, કારણ કે તે માપી ન શકાય
એવા ગણની વિટાલીની \citeyear{Vitali1905}ની
રચનામાંથી મળે છે. જોકે વિટાલી અને
બાનાખ--ટાર્સ્કીના પુરાવાના ગણસિદ્ધાંતના
પાસાં ખૂબ સરખાં છે. બંને પરિણામો
વચ્ચેનો મુખ્ય ભેદ એટલો જ છે કે
બાનાખ--ટાર્સ્કી \emph{સાન્ત} વિભાજન
લે છે, જ્યારે વિટાલીનો વિરોધાભાસ
\emph{ગણનીય અનંત} વિભાજન લે છે.
\citet{Weston2003} કહે છે તેમ,
વિટાલીનો વિરોધાભાસ ``બાનાખ--ટાર્સ્કી
જેટલો આઘાતજનક તો ચોક્કસ નથી,
પણ તે બતાવે છે કે ભૌમિતિક
વિરોધાભાસો `સરળ' પરિસ્થિતિઓમાં
પણ થઈ શકે છે.''

વિટાલીનો વિરોધાભાસ દ્વિપરિમાણીય
આકૃતિ, વર્તુળ, અંગે છે. તેથી
સમવતલ $\Real^2$માં કામ કરીએ.
$\rotationsgroup$ને મૂળબિંદુની આસપાસ
બિંદુઓના ઘડિયાળની દિશામાં એવા
પરિભ્રમણોનો ગણ લો, જેમના કોણ
$[0,2\pi)$માં હોય અને પૂર્ણ
પરિભ્રમણના પરિમેય ભાગ હોય.
$\rotationsgroup$ વિશેનાં કેટલાક
બીજગણિતીય તથ્યો આ છે (વિધાન
ન સમજાય તો પુરાવાથી તેનો અર્થ
સ્પષ્ટ થશે):
\sourcecorrection{OLMD-151}{મૂળમાં
``$[0,2\pi)$ વચ્ચેના પરિમેય રેડિયન''
કહ્યાં છે. એ પરિભ્રમણો સમૂહ નથી:
પરિમેય $r$નો પ્રતિપરિભ્રમણ
$2\pi-r$ સામાન્ય રીતે પરિમેય
નથી. અહીં પૂર્ણ પરિભ્રમણના પરિમેય
ભાગ $2\pi q$ લીધા છે; શૂન્ય કોણનો
પ્રતિપરિભ્રમણ પણ શૂન્ય જ છે.}

\begin{lem}\ollabel{rotationsgroupabelian}
વિધેયોના સંયોજન હેઠળ
$\rotationsgroup$ એક ક્રમવિનિમેય
{સમૂહ} છે.
\end{lem}

\begin{proof}
$0$ રેડિયનના પરિભ્રમણને
$0_{\rotationsgroup}$ લખીએ.
તે $\rotationsgroup$નો એકમ ઘટક છે,
કારણ કે કોઈ પણ $\rho \in \rotationsgroup$
માટે $\comp{0_{\rotationsgroup}}{\rho} =
\comp{\rho}{0_{\rotationsgroup}} =
\rho$.

દરેક ઘટકનું પ્રતિઘટક છે. જો
$\rho \in \rotationsgroup$ એ $r$
રેડિયનથી ફેરવે, તો કોણ શૂન્ય ન હોય ત્યારે
$\rho^{-1} \in \rotationsgroup$ એ
$2\pi-r$ રેડિયનથી ફેરવે;
શૂન્ય કોણ માટે તે શૂન્યથી જ ફેરવે.
તેથી $\rho \circ \rho^{-1} =
0_\rotationsgroup$.

સંયોજન સાહચર્ય ધરાવે છે: કોઈ પણ
$\rho, \sigma, \tau \in
\rotationsgroup$ માટે
$\comp{\rho}{(\comp{\sigma}{\tau})} =
\comp{(\comp{\rho}{\sigma})}{\tau}$.

સંયોજન ક્રમવિનિમેય છે: કોઈ પણ
$\rho, \sigma \in \rotationsgroup$
માટે $\comp{\rho}{\sigma} =
\comp{\sigma}{\rho}$.
\end{proof}

ખરેખર આપણા સમૂહ
$\rotationsgroup$ને બે ભાગમાં વહેંચીને,
બેમાંથી કોઈ પણ ભાગ વડે આખો
સમૂહ પાછો મેળવી શકીએ:

\begin{lem}\ollabel{disjointgroup}
$\rotationsgroup$નું બે વિચ્છિન્ન ગણો
$\rotationsgroup_{1}$ અને
$\rotationsgroup_{2}$માં વિભાજન છે;
બંને આખો $\rotationsgroup$
પાછો મેળવવાનો આધાર આપે છે.
\end{lem}

\begin{proof}
$\rotationsgroup_{1}$માં
$[0, \pi)$માં આવેલા પૂર્ણ
પરિભ્રમણના પરિમેય ભાગના
કોણવાળાં પરિભ્રમણો લો;
$\rotationsgroup_{2} = \rotationsgroup
\setminus \rotationsgroup_1$ લો.
પ્રાથમિક બીજગણિતથી
$\Setabs{\comp{\rho}{\rho}}{\rho \in \rotationsgroup_1} =
\rotationsgroup$. એ જ પરિણામ
$\rotationsgroup_2$ માટે મળે છે.
\end{proof}

સમૂહો અંગેનું આ તથ્ય
\olref{vitaliparadox} સાબિત કરવા
વાપરીશું. $\onesphere$ને એકમ
વર્તુળ કહીએ, એટલે સમવતલના મૂળબિંદુથી
બરાબર~$1$ અંતરે આવેલા બિંદુઓનો
ગણ, એટલે
$\Setabs{\tuple{r,s} \in \Real^2}{\sqrt{r^2+s^2}=1}$.
$\onesphere$ પરનો નીચેનો સંબંધ
વિચારી $\onesphere$ને ભાગોમાં વહેંચીશું:
\[
	r \sim s \emph{ ત્યારે અને ત્યારે જ }(\exists \rho \in \rotationsgroup)\rho(r) = s.
\]
એટલે $\onesphere$ના બે બિંદુઓ
આ સંબંધથી જોડાયેલા હોય છે ત્યારે
અને ત્યારે જ, જ્યારે મૂળબિંદુની આસપાસ
પરિમેય ભાગના પૂર્ણ પરિભ્રમણથી
એકમાંથી બીજું મળી શકે. આશ્ચર્ય નહીં:

\begin{lem}
$\sim$ સમકક્ષતા સંબંધ છે.
\end{lem}

\begin{proof}
\olref{rotationsgroupabelian}થી તરત મળે છે.
\end{proof}

હવે પસંદગી વાપરી $\onesphere$ના
$\sim$ હેઠળના દરેક સમકક્ષતા વર્ગમાંથી
બરાબર એક ઘટક ધરાવતો ગણ $C$ લઈએ.
એટલે સમકક્ષતા વર્ગોના
ગણ\footnote{$\rotationsgroup$ !!{enumerable} છે,
તેથી $E$નો દરેક !!{element}
!!{enumerable} છે.
$\onesphere$ !!{nonenumerable} હોવાથી
\olref{vitalicover} અને
\olref[card-arithmetic][simp]{kappaunionkappasize}
મુજબ $E$ !!{nonenumerable} છે. તેથી અહીં
\emph{અ}ગણનીય પસંદગી વપરાય છે.}
\[
	E = \Setabs{\equivrep{r}{\sim}}{r \in \onesphere},
\]
પર પસંદગી વિધેય $f$ લઈ
$C = \ran{f}$ મૂકીએ.
દરેક $\rho \in \rotationsgroup$
માટે $\funimage{\rho}{C}$ એ $C$ના
દરેક બિંદુ પર $\rho$ લાગુ કરતાં
મળેલા બિંદુઓનો ગણ છે. આગળનાં
બે પરિણામો બતાવે છે કે આ ગણો
આખું વર્તુળ ઢાંકે છે અને તેમનો
આપસમાં છેદ નથી:

\begin{lem}\ollabel{vitalicover}
$\onesphere = \bigcup_{\rho \in \rotationsgroup} \funimage{\rho}{C}$.
\end{lem}

\begin{proof}
$s \in \onesphere$ નિશ્ચિત લો;
કોઈ $r \in C$ એવું છે કે
$r \in \equivrep{s}{\sim}$,
એટલે $r \sim s$, એટલે કોઈ
$\rho \in \rotationsgroup$ માટે
$\rho(r) = s$.
\end{proof}

\begin{lem}\ollabel{vitalinooverlap}
$\rho_1 \neq \rho_2$ હોય, તો
$\funimage{\rho_1}{C} \cap \funimage{\rho_2}{C} = \emptyset$.
\end{lem}

\begin{proof}
$s \in \funimage{\rho_{1}}{C} \cap
\funimage{\rho_{2}}{C}$ ધારો.
તેથી કોઈ $r_{1}, r_{2} \in C$ માટે
$s = \rho_{1}(r_{1}) = \rho_{2}(r_{2})$.
આથી $\rho^{-1}_2(\rho_1(r_1)) = r_2$
અને $\comp{\rho_1}{\rho^{-1}_2} \in
\rotationsgroup$, એટલે $r_{1} \sim
r_{2}$. $C$ એ $\sim$ હેઠળના
દરેક સમકક્ષતા વર્ગમાંથી
બરાબર એક ઘટક લે છે, એટલે
$r_1 = r_2$. તેથી
$s = \rho_1(r_1) = \rho_2(r_1)$,
અને પરિણામે $\rho_1 = \rho_2$.
\end{proof}

હવે અગાઉનાં બીજગણિતીય તથ્યો
વર્તુળ પર લાગુ કરીએ:

\begin{lem}\ollabel{pseudobanachtarski}
$\onesphere$નું બે વિચ્છિન્ન ગણો
$D_{1}$ અને $D_{2}$માં એવું
વિભાજન છે કે $D_{1}$ને ગણનીય
ગણોમાં વહેંચી, તેમને ફેરવીને
$\onesphere$ની એક નકલ બનાવી
શકાય છે (અને $D_{2}$ માટે પણ
એવું જ થાય છે).
\end{lem}

\begin{proof}
\olref{disjointgroup}ના
$\rotationsgroup_{1}$ અને
$\rotationsgroup_{2}$ વાપરી આ લો:
\begin{align*}
	D_{1} &= \bigcup_{\rho \in \rotationsgroup_1} \funimage{\rho}{C} & 
	D_{2} &= \bigcup_{\rho \in \rotationsgroup_2} \funimage{\rho}{C}
\end{align*}
\olref{vitalicover}થી આ
$\onesphere$નું વિભાજન છે, અને
\olref{vitalinooverlap}થી
$D_1$ અને $D_2$ વિચ્છિન્ન છે.
રચના પ્રમાણે $D_1$ને દરેક
$\rho \in \rotationsgroup_1$ માટેના
$\funimage{\rho}{C}$ જેવા ગણનીય
ઘણા ગણોમાં વહેંચી શકાય છે.
તેમને ફેરવી $\onesphere$ની
નકલ બનાવી શકાય છે, કારણ કે
\olref{disjointgroup} અને
\olref{vitalicover}થી
$\onesphere = \bigcup_{\rho \in
\rotationsgroup}\funimage{\rho}{C} =
\bigcup_{\rho \in
\rotationsgroup_1}\funimage{(\comp{\rho}{\rho})}{C}$.
$D_2$ માટે પણ એ જ દલીલ છે.
\end{proof}\noindent
આમાંથી વિટાલીનો વિરોધાભાસ
તરત મળે છે. $\onesphere$ને
ગણનીય ટુકડાઓમાં (જુદા જુદા
$\funimage{\rho}{C}$માં) વહેંચીને
પછી તેમને અચળ આકારવાળી ગતિથી
ખસેડીએ: $D_1$નો દરેક ટુકડો
ફેરવીએ અને $D_2$નો દરેક ટુકડો
પહેલાં સ્થાનાંતરિત કરી પછી
ફેરવીએ. આમ $\onesphere$માંથી
તેની \emph{બે} નકલો મળે છે.
\sourcecorrection{OLMD-152}{મૂળ પુરાવામાં
ટુકડાઓના સૂચક માટે $R_1$ લખ્યું
છે, પણ એવો ગણ વ્યાખ્યાયિત થયો
નથી. અહીં અગાઉનો
$\rotationsgroup_1$ જ વાપર્યો છે.}

પુરાવાની વ્યૂહરચના ફરી જોઈએ.
આપણે સમવતલના પરિભ્રમણ-સમૂહ
વિશે કેટલાંક બીજગણિતીય તથ્યોથી
શરૂઆત કરી. એ સમૂહથી
$\onesphere$ને સમકક્ષતા વર્ગોમાં
વહેંચ્યું. પછી દરેક વર્ગમાંથી
ઘટક પસંદ કરવા પસંદગી વાપરીને
``વિરોધાભાસ'' સુધી પહોંચ્યા.

બાનાખ--ટાર્સ્કી સાબિત કરવા પણ
બરાબર એ જ વ્યૂહરચના વાપરીએ.
મુખ્ય ભેદ એ છે કે બાનાખ--ટાર્સ્કી
માટેનાં બીજગણિતીય તથ્યો
વિટાલીના વિરોધાભાસ માટેનાં
તથ્યો કરતાં ઘણાં જટિલ છે.
પણ એ તથ્યોનો પસંદગી સાથે
કોઈ સંબંધ નથી. તેમને ટૂંકમાં
જણાવીએ.

બાનાખ--ટાર્સ્કી સાબિત કરવા
સૌ પહેલાં \olref{disjointgroup}ને
સમાન પરિણામ સ્થાપીએ: કોઈ પણ
\emph{મુક્ત સમૂહ}ને ચાર ભાગમાં
વહેંચી શકાય છે, જેમને સહજ રીતે
``ખસેડીને'' આખા સમૂહની બે
નકલો પાછી મેળવી શકીએ.\footnote{
\emph{ચાર} ભાગ પૂરતા છે એ
\cite{Robinson1947}નું પરિણામ છે.
તાજેતરના પુરાવા માટે
\citet[Theorem 5.2]{Wagon2016} જુઓ.
સમૂહના ભાગોને ``ખસેડવા''
એ રીતે વર્ણવવામાં આપણે
\citet[p.~3]{Weston2003}ને
અનુસરીએ છીએ.}
પછી બતાવીએ કે $\Real^3$ના
મૂળબિંદુની આસપાસનાં બે ચોક્કસ
પરિભ્રમણોથી પરિભ્રમણોનો
મુક્ત સમૂહ $F$ બને છે.\footnote{જુઓ
\citet[Theorem 2.1]{Wagon2016}.}
(અહીં સુધી પસંદગી વાપરી નથી.)
હવે ગોળાની સપાટી પરના બે
બિંદુઓને ``સમાન'' કહીએ ત્યારે
અને ત્યારે જ, જ્યારે $F$નું
કોઈ પરિભ્રમણ એકને બીજામાં લઈ
જાય. પછી સમાન બિંદુઓના
દરેક સમકક્ષતા વર્ગમાંથી બરાબર
એક બિંદુ પસંદ કરવા
\emph{પસંદગી વાપરીએ}. $F$ના
વિભાજનને ગોળાની સપાટી પર
\olref{pseudobanachtarski}ની
જેમ લાગુ કરતાં સપાટી ચાર
ભાગમાં વહેંચાય છે; તેમને
``ખસેડીને'' સપાટીની બે નકલો
મેળવીએ. અને આમ
\citep{Hausdorff1914}નું પરિણામ
સ્થાપિત થાય છે:
\sourcecorrection{OLMD-153}{આ ટૂંકા
રેખાચિત્રમાં એક આવશ્યક પગલું
છૂટી ગયું છે: $F$ના અતુચ્છ
પરિભ્રમણોથી સ્થિર રહેતા બિંદુઓ
પર સમૂહની ક્રિયા મુક્ત નથી.
એવા ગણનીય બિંદુઓ પહેલાં
અલગ કરવાના અને અંતે પાછા
સમાવવાના તર્ક વિના માત્ર
દરેક વર્ગમાંથી એક પ્રતિનિધિ
લેવાથી આખી સપાટીનું નીચેનું
પરિણામ સાબિત થતું નથી.}

\begin{thm}[હાઉસડૉર્ફનો વિરોધાભાસ ($\ZFC$માં)]
કોઈ પણ ગોળાની સપાટીને
સાન્ત સંખ્યામાં ટુકડાઓમાં
વિભાજિત કરી શકાય છે; એ
ટુકડાઓને (પરિભ્રમણ અને
સ્થાનાંતરણથી) ફરી ગોઠવી
તે ગોળાની બે વિચ્છિન્ન
નકલો બનાવી શકાય છે.
\end{thm}

સંપૂર્ણ બાનાખ--ટાર્સ્કી પ્રમેય
મેળવવા થોડા વધુ બીજગણિતીય
ઉપાયો જોઈએ (તે માત્ર સપાટી
જ નહીં, અંદરનો ભાગ પણ લે છે).
ખુલ્લું કહીએ તો એ બીજગણિતીય
કેક પરની સજાવટ જેવી બાબત છે.
તેથી વેસ્ટન લખે છે:
\begin{quote}
  [\ldots] મુક્ત સમૂહો અંગેનું
  પરિણામ બાનાખ--ટાર્સ્કીના
  વિરોધાભાસના પુરાવાનું
  \emph{મુખ્ય પગલું} છે. આ
  દૃષ્ટિએ બાનાખ--ટાર્સ્કીનો
  વિરોધાભાસ $\Real^3$ વિશેના
  વિધાન કરતાં [$\Real^3$માં
  સ્થાનાંતરણ અને પરિભ્રમણોના]
  સમૂહની જટિલતા વિશેનું
  વિધાન વધુ છે.
  \cite[p.~16]{Weston2003}
\end{quote}
એટલે \emph{સાન્ત} વિભાજન
(બાનાખ--ટાર્સ્કીની જેમ) મળે
કે \emph{ગણનીય અનંત} વિભાજન
(વિટાલીના વિરોધાભાસની જેમ),
તે દ્વિપરિમાણીય કે ત્રિપરિમાણીય
અવકાશમાં કામ કરતાં સમૂહસિદ્ધાંતના
અમુક તથ્યો પર આધાર રાખે છે.

સ્વીકારીએ કે આ છેલ્લું નિરીક્ષણ
\olref[banach]{sec}ના અંતની
રમૂજ થોડી બગાડે છે. દ્વિપરિમાણીય
હોવાથી ``Banach-Tarski''ને
``Banach-Tarski Banach-Tarski''
બનાવવા તેના \emph{ગણનીય અનંત}
ટુકડા કરવા પડે. રમૂજ સુધારવા
ત્રિપરિમાણમાં લખવું પડે.
આ કસરત વાચક માટે છોડીએ.

એક છેલ્લી ટિપ્પણી.
\olref[banach]{sec}માં આપણે કહ્યું
હતું કે ગોળાના મળતા ``ટુકડાઓ''
\emph{માપનીય} હોઈ શકતા નથી;
તે ચિત્રમાં ન ઉતારી શકાય એવા
``અનંત વિખેરાયેલા'' ગણો છે.
\olref{pseudobanachtarski}
મેળવવામાં પસંદગીના આપણા
ઉપયોગ માટે પણ એવું જ છે.
આ બાબત સમજાવવા જેવી છે.

ફરી થોડું પૃષ્ઠભૂમિ-જ્ઞાન
આપવું પડે (પરંતુ આ
\emph{માત્ર} રૂપરેખા છે; તમે
\emph{માપ} વિશેનો પાઠ્યપુસ્તકનો
વિભાગ જોઈ શકો). ગણ $X$ માટે
માપ વ્યાખ્યાયિત કરવું એટલે
$X$ પરના કોઈ ``$\sigma$-બીજગણિત''ના
દરેક $E$ને અઋણ મૂલ્ય
$\mu(E) \in \Real$ આપવું.
વિગતો અહીં જરૂરી નથી;
ફક્ત વિધેય $\mu$એ ગણનીય
યોગશીલતા પાળવી પડે: પરસ્પર
વિચ્છિન્ન ગણોના ગણનીય યોગગણનું
માપ તેમના અલગ માપોનો સરવાળો,
એટલે $X_n$ વિચ્છિન્ન હોય ત્યારે
$\mu(\bigcup_{n < \omega} X_n) =
\sum_{n < \omega}\mu(X_n)$.
કોઈ ગણ ``અમાપનીય'' છે એટલે
તેને યોગ્ય રીતે માપ આપી
શકાતું નથી. હવે અગાઉનો
$\rotationsgroup$ વાપરીએ:

\begin{cor}[વિટાલી]
$\mu(\onesphere) = 1$ એવું
માપ $\mu$ લો, જેમાં $X$ અને
$Y$ સમરૂપ હોય ત્યારે
$\mu(X) = \mu(Y)$. તો દરેક
$\rho \in \rotationsgroup$ માટે
$\funimage{\rho}{C}$ અમાપનીય છે.
\end{cor}

\begin{proof}
વિરોધાભાસ માટે ઊલટું ધારો.
એટલે કોઈ $\sigma \in
\rotationsgroup$ અને કોઈ
$r \in \Real$ માટે
$\mu(\funimage{\sigma}{C}) = r$
લો. કોઈ પણ $\rho \in
\rotationsgroup$ માટે
$\funimage{\rho}{C}$ અને
$\funimage{\sigma}{C}$ સમરૂપ
છે; તેથી દરેક $\rho \in
\rotationsgroup$ માટે
$\mu(\funimage{\rho}{C}) = r$.
\olref{vitalicover} અને
\olref{vitalinooverlap} મુજબ
$\onesphere = \bigcup_{\rho \in
\rotationsgroup}\funimage{\rho}{C}$
પરસ્પર વિચ્છિન્ન ગણોનો
ગણનીય યોગગણ છે. તેથી
ગણનીય યોગશીલતાથી
$\mu(\onesphere) = 1$ એ
દરેક $\funimage{\rho}{C}$ના
માપોનો સરવાળો છે, એટલે
\[
	1 = \mu(\onesphere) = \sum_{\rho \in \rotationsgroup}\mu(\funimage{\rho}{C}) = \sum_{\rho \in \rotationsgroup}r
\]
પણ $r = 0$ હોય તો
$\sum_{\rho \in \rotationsgroup}r = 0$,
અને $r > 0$ હોય તો
$\sum_{\rho \in \rotationsgroup}r = \infty$.
\sourcecorrection{OLMD-154}{મૂળ પુરાવામાં
બે વાર $\rho\in C$ લખાયું છે,
પણ $C$માં બિંદુઓ છે અને
$\rho$ પરિભ્રમણ હોવું જોઈએ;
અહીં બંને જગ્યાએ
$\rho\in\rotationsgroup$ કર્યું
છે. માપના મૂલ્યો અઋણ
હોવાની માનક શરત પણ સ્પષ્ટ
કરી છે, જેથી અંતિમ $r=0$
અથવા $r>0$ વિભાજન પૂર્ણ છે.}
\end{proof}

\end{document}
